
\begin{exercise}[Counterexample joint convergence \Coffeecup]
    \label{ex: counterexample joint convergence}
    Let \(X\) be a random variable with \(\prob{X=-1} = \prob{X=1} = \frac12\)
    and define \(X_n \coloneq X\) and \(Y_n \coloneq -X\). Show that \(X_n
    \overset{d}\to X\) and
    \(Y_n \overset{d}\to X\) but 
    \[
        X_n + Y_n \overset{d}{\not\to} X + X = 2X.
    \]
    Conclude that \((X_n, Y_n)\) does not converge to \((X,X)\) in distribution.
    \begin{remark}
        Observe that \(X_n \overset{p}\to X\). It is therefore not sufficient
        that one of the sequences converges in probability as Slutzky (Theorem
        \ref{thm: joint convergence}~\ref{it: joint conv law slutsky}) may
        suggest.
    \end{remark}
\end{exercise}

\begin{solution}
    Clearly, \(X_n + Y_n = 0\) for all \(n\). In particular we have
    \[
        \prob{X_n + Y_n \in A} = \ind{A}(0)
        \neq \frac12\ind{A}(2) + \frac12 \ind{A}(-2)
        = \prob{2X \in A}.
    \]
    for \(A= \{0\}\) for example. Since \(\prob{2X = 0} = 0\) this is
    however a requirement for \(X_n + Y_n \overset{d}\to 2X\) by \ref{it:
    s.zero boundary}. Thus \(X_n + Y_n \overset{d}{\not\to} 2X\).  This also
    implies that \((X_n, Y_n)\) does not converge to \((X,X)\) in
    distribution as otherwise we would have \(X_n + Y_n \overset{d}\to X +
    X\) by the continuous mapping theorem.
\end{solution}
